% БҮЛЭГ 3 — СИСТЕМИЙН АРХИТЕКТУР БА ЗАГВАРЫН ЗОХИОН БАЙГУУЛАЛТ

\chapter{Системийн архитектур ба загварын зохион байгуулалт}
\label{Chapter3}

Энэхүү бүлэгт судалгааны ажлын хүрээнд боловсруулсан туршилтын архитектур, загварын
сургалтын дамжлага, интерфейсийн бүтэц, болон үнэлгээний арга зүйг нарийвчлан
тайлбарлана. Судалгааны үндсэн таамаглал (hypothesis) нь хоёр шатлалт ангиллын
архитектур ашиглан ``Сөрөг'' агуулгыг нарийвчлан ялгах замаар системийн
ерөнхий нарийвчлалыг нэмэгдүүлэх байсан бөгөөд уг бүтцийг дараах байдлаар
зохион байгуулсан.

\textbf{Чухал тэмдэглэл (архитектурын сонголт).} Энэ бүлэгт танилцуулж буй
хоёр шатлалт каскад архитектур нь судалгааны \emph{анхны санал болгосон}
шийдэл бөгөөд түүнийг нэг шатлалт (Flat) хувилбартай туршилтаар харьцуулсан.
Хоёр шатлалт хувилбар нь каскад алдааны дамжуулалтын улмаас нийт гүйцэтгэл
$78.61\%$ болж буурсан тул \textbf{эцсийн шийдэл болгон нэг шатлалт MN-BERT
загварыг сонгосон} (үндэслэл, тоон харьцуулалтыг бүлэг~\ref{Chapter5} ба
бүлэг~\ref{Chapter6}-д дэлгэрэнгүй үзүүлэв). Иймд энэ бүлгийн хоёр шатлалт
тайлбар нь харьцуулалтын суурь болохоос биш эцсийн системийн архитектур биш юм.

%-------------------------------------------------------------------------------
\section{Системийн туршилтын архитектур}

Системийн зохион байгуулалтыг илүү тодорхой болгохын тулд эхлээд системийн
хэрэглээний тохиолдлын (Use Case) болон үйл ажиллагааны (Activity) логикийг
тодорхойлов.

\subsection{Хэрэглээний тохиолдлын диаграмм (Use Case Diagram)}

Зураг \ref{fig:use-case}-т харуулсанчлан систем нь гурван үндсэн оролцогчтой:
Өсвөр насны хэрэглэгч, Модератор (зохицуулагч) болон Системийн администратор.
Хэрэглэгч сэтгэгдэл оруулах бол, Модератор нь системийн ангилсан үр дүнг хянаж,
хортой агуулгыг шүүх, бүтээлч шүүмжлэлийг зөвшөөрөх үйлдлийг гүйцэтгэнэ.

\begin{figure}[H]
\centering
\begin{tikzpicture}[
  usecase/.style={ellipse, draw, minimum width=5.0cm, minimum height=1.0cm,
                  fill=blue!8, draw=blue!50, line width=0.8pt,
                  align=center, font=\small},
  actor/.style={circle, draw, minimum size=0.7cm,
                fill=orange!14, draw=orange!65, line width=0.8pt},
  arr/.style={draw, ->, >=Stealth, thick, line width=0.7pt},
]

  % ── SYSTEM BOUNDARY ──────────────────────────────────────────────────────
  \draw[rounded corners=6pt, draw=gray!55, fill=gray!4, line width=1pt]
    (-3.2, -4.6) rectangle (3.2, 3.6);
  \node[font=\small\bfseries, text=gray!65] at (0, 3.25)
    {NLP Ангиллын Систем};

  % ── USE CASES ────────────────────────────────────────────────────────────
  \node[usecase] (uc1) at (0,  2.2) {Сэтгэгдэл оруулах};
  \node[usecase] (uc2) at (0,  0.8) {Ангиллын дүн үзэх};
  \node[usecase] (uc3) at (0, -0.7) {Хортой агуулга шүүх};
  \node[usecase] (uc4) at (0, -2.2) {Бүтээлч шүүмжлэл зөвшөөрөх};
  \node[usecase] (uc5) at (0, -3.8) {Загвар сургах / шинэчлэх};

  % ── LEFT ACTORS ──────────────────────────────────────────────────────────
  \node[actor] (user)  at (-6.0,  2.2) {};
  \node[font=\small, align=center] at (-6.0, 1.35)
    {Өсвөр насны\\Хэрэглэгч};

  \node[actor] (admin) at (-6.0, -3.8) {};
  \node[font=\small, align=center] at (-6.0, -4.65)
    {Системийн\\Администратор};

  % ── RIGHT ACTOR ──────────────────────────────────────────────────────────
  \node[actor] (mod)   at ( 6.0, -0.7) {};
  \node[font=\small, align=center] at ( 6.0, -1.55)
    {Модератор};

  % ── CONNECTIONS ──────────────────────────────────────────────────────────
  \draw[arr] (user)  -- (uc1.west);
  \draw[arr] (admin) -- (uc5.west);
  \draw[arr] (mod)   -- (uc2.east);
  \draw[arr] (mod)   -- (uc3.east);
  \draw[arr] (mod)   -- (uc4.east);

\end{tikzpicture}
\caption{Системийн хэрэглээний тохиолдлын диаграмм: Хэрэглэгч сэтгэгдэл оруулж,
  Модератор ангиллын үр дүнг хянан шүүж, Администратор загварыг удирдана.}
\label{fig:use-case}
\end{figure}

\subsection{Үйл ажиллагааны диаграмм (Activity Diagram)}

Системийн логик урсгалыг Зураг \ref{fig:activity-diagram}-т харуулав. Энэхүү
диаграмм нь хоёр шатлалт ангиллын шийдвэр гаргах үйл явцыг нарийвчлан харуулж
байна.

\begin{figure}[H]
    \centering
    \resizebox{\linewidth}{!}{%
    \begin{tikzpicture}[
        node distance=1cm and 1.5cm,
        start/.style={circle, draw, fill=black, inner sep=3pt},
        stop/.style={circle, draw, double, fill=black!50, inner sep=2.5pt},
        procstep/.style={rectangle, draw, rounded corners, minimum width=4cm, minimum height=1cm, text centered, fill=green!5, drop shadow, align=center},
        decision/.style={diamond, draw, aspect=2, minimum width=3cm, text centered, fill=yellow!10, font=\small, drop shadow, inner sep=0pt, align=center},
        arrow/.style={thick, ->, >=stealth}
    ]
        \node (init) [start] {};
        \node (input) [procstep, below=of init] {Текст хүлээн авах};
        \node (prep) [procstep, below=of input] {Урьдчилсан боловсруулалт};
        \node (s1) [decision, below=1.5cm of prep] {Сөрөг? (Stage 1)};
        
        \node (allow) [procstep, left=2.5cm of s1] {Зөвшөөрөх (Allow)};
        
        \node (s2) [decision, below=1.5cm of s1] {Хортой? (Stage 2)};
        \node (block) [procstep, right=2.5cm of s2] {Блок хийх (Block)};
        
        \node (end) [stop, below=of s2] {};

        \draw [arrow] (init) -- (input);
        \draw [arrow] (input) -- (prep);
        \draw [arrow] (prep) -- (s1);
        \draw [arrow] (s1) -- node[anchor=east, xshift=-0.1cm] {Тийм} (s2);
        \draw [arrow] (s1) -- node[anchor=south] {Үгүй} (allow);
        \draw [arrow] (s2.west) -| node[anchor=south, xshift=1.5cm] {Үгүй} (allow);
        \draw [arrow] (s2) -- node[anchor=south] {Тийм} (block);
        
        \draw [arrow] (allow) |- (end);
        \draw [arrow] (block) |- (end);
    \end{tikzpicture}}
    \caption{Системийн шүүлтүүрийн логикийн үйл ажиллагааны диаграмм}
    \label{fig:activity-diagram}
\end{figure}

\subsection{Давхаргын архитектур}

Системийн архитектур нь гурван үндсэн давхаргаас бүрдэнэ: (1)~өгөгдлийн
боловсруулалтын давхарга (Data Processing Layer), (2)~загварын давхарга (Model
Layer), (3)~хэрэглэгчийн интерфейсийн давхарга (Presentation Layer). Эдгээр
давхарга нь бие биенээс харьцангуй бие даасан (loosely coupled) бүтэцтэй бөгөөд
давхарга бүрийг тус тусад нь шинэчлэх, туршихад хялбар.

\begin{figure}[H]
  \centering
  \begin{tikzpicture}[
      block/.style={rectangle, draw, fill=blue!10, text width=10em,
                    text centered, rounded corners, minimum height=3em, align=center},
      line/.style={draw, -latex}
    ]
    \node [block] (input) {Хам текст \\ (найман эх сурвалж)};
    \node [block, below of=input, node distance=3cm] (preprocess)
          {Урьдчилсан боловсруулалт \\ (Цэвэрлэгээ, Токенчлол)};
    \node [block, below of=preprocess, node distance=3cm] (stage1)
          {Stage~1: Сентимент ангилал \\ (Эерэг / Саармаг / Сөрөг)};
    \node [block, below of=stage1, node distance=3cm] (stage2)
          {Stage~2: Нарийвчилсан ангилал \\ (Хортой сөрөг vs Бүтээлч шүүмжлэл)};
    \node [block, below of=stage2, node distance=3cm] (output)
          {Үр дүн \\ (Gradio интерфейс)};
    \path [line] (input) -- (preprocess);
    \path [line] (preprocess) -- (stage1);
    \path [line] (stage1) -- node [right, xshift=0.3cm]
          {\small Сөрөг} (stage2);
    \path [line] (stage2) -- (output);
  \end{tikzpicture}
  \caption{Системийн ерөнхий архитектурын схем}
  \label{fig:system-arch}
\end{figure}

Зураг~\ref{fig:pipeline}-д хам текстээс эцсийн шошго хүртэлх нарийвчилсан
дамжлагыг харуулав.

\begin{figure}[H]
\centering
\begin{tikzpicture}[
  node distance=0.5cm,
  box/.style={draw, rounded corners, minimum width=4.8cm, minimum height=0.75cm,
              text centered, font=\small, fill=blue!8},
  wide/.style={draw, rounded corners, minimum width=4.8cm, minimum height=0.75cm,
               text centered, font=\small, fill=green!12},
  arrow/.style={->, >=Stealth, thick}
]
\node[box]  (raw)   at (0,0)    {Хам текст (найман эх сурвалж)};
\node[box]  (filt)  at (0,-1.3) {Кирилл шүүлт ($\geq$85\% кирилл)};
\node[box]  (clean) at (0,-2.6) {Цэвэрлэгээ (URL, HTML, emoji → placeholder)};
\node[box]  (tok)   at (0,-3.9) {SentencePiece токенчлол (MN-BERT)};
\node[box]  (enc)   at (0,-5.2) {\texttt{[CLS]}+токен дараалал → 12-давхарга энкодер};
\node[wide] (cls)   at (0,-6.5) {\texttt{[CLS]} вектор ($\mathbb{R}^{768}$) → Linear + Softmax};
\node[wide] (out)   at (0,-7.8) {Ангиллын шошго (Эерэг / Саармаг / Хортой сөрөг / Бүтээлч шүүмжлэл)};
\draw[arrow] (raw)   -- (filt);
\draw[arrow] (filt)  -- (clean);
\draw[arrow] (clean) -- (tok);
\draw[arrow] (tok)   -- (enc);
\draw[arrow] (enc)   -- (cls);
\draw[arrow] (cls)   -- (out);
\end{tikzpicture}
\caption{MN-BERT сургалтын pipeline (хам текстээс ангиллын шошго хүртэл)}
\label{fig:pipeline}
\end{figure}

Зураг~\ref{fig:system-arch}-т үзүүлсэн ерөнхий бүтцийн дагуу хэрэглэгчийн
оруулсан текст эхлээд урьдчилсан боловсруулалтын шатыг давж, дараа нь хоёр
дараалсан ангиллын шатаар дамжина. Stage~1 нь текстийг Эерэг, Саармаг,
Сөрөг гэсэн гурван ангилалд хуваана. Зөвхөн Сөрөг гэж ангилагдсан
текст Stage~2 руу дамжих бөгөөд тэнд Хортой сөрөг (хортой) ба Бүтээлч шүүмжлэл
(бүтээлч шүүмжлэл) гэсэн хоёр ангилалд ялгана.

%-------------------------------------------------------------------------------
\section{Хоёр шатлалт ангиллын архитектур}

Хоёр шатлалт архитектурыг сонгосон гол шалтгаан нь дөрвөн ангиллын хоорондын
семантик ялгааны тэнцвэргүй байдал юм. Нэг загвараар бүгдийг ялгахад
Эерэг болон Саармаг нь Сөрөг-ээс харьцангуй амархан ялгагддаг бол Хортой сөрөг болон
Бүтээлч шүүмжлэл хоёрыг бие биенээс ялгах нь хэл шинжлэлийн
хувьд ихэд төвөгтэй. Учир нь хоёр ангилал хоёулаа сөрөг өнгө аястай үгсийг
агуулдаг боловч гол ялгаа нь илтгэгчийн зорилгод оршдог: Хортой сөрөг нь
хохирол учруулах, доромжлох зорилготой бол Бүтээлч шүүмжлэл нь сайжруулах,
тодорхой асуудлыг шийдвэрлэхэд чиглэсэн байдаг. Гадаргуугийн лексик
онцлогуудаар --- тухайлбал хэрэглэгдэж буй үгсийн утга --- энэ хоёрыг найдвартай
ялгах боломжгүй тул контекст мэдрэмжтэй загвар зайлшгүй шаардлагатай.
Тиймд дамжлагыг хоёр тусдаа шатанд хуваах нь загвар бүрийг өөрийн тусгай
даалгаварт төвлөрүүлэх боломж олгоно.

\subsection{Stage~1 --- Сентимент ангилал}

Stage~1 загварын зорилго нь оролтын текстийг Эерэг, Саармаг, Сөрөг гэсэн
гурван ангилалд хуваах юм. Энэ шатанд дараах загваруудыг хэрэгжүүлж,
харьцуулна:

\begin{itemize}
  \item \textbf{MN-BERT (нарийн тохируулгатай):} Монгол хэлний урьдчилсан сургалттай
        (pre-trained) BERT загварыг гурван ангиллын өгөгдлийн багц дээр нарийн тохируулга
        хийнэ. Сүүлийн \texttt{[CLS]} токены давхаргын гаралтыг softmax
        ангилагч руу дамжуулна.
  \item \textbf{Logistic Regression (LR):} TF-IDF шинж чанарын дээр L2
        регуляризацитай LR загвар сургана.
  \item \textbf{Naive Bayes (NB):} TF-IDF шинж чанарын дээр Multinomial
        Naive Bayes ашиглана.
  \item \textbf{SVM:} Linear kernel бүхий Support Vector Machine загвар
        сургана.
\end{itemize}

Stage~1-ийн гол шаардлага нь Сөрөг ангилалын recall-ийг аль болох
өндөр байлгах явдал юм. Сөрөг текстийг алдвал (False Negative)
тухайн текст Stage~2 руу дамжихгүй, улмаар Хортой сөрөг агуулга илрүүлэгдэхгүй болно.

\subsection{Stage~2 --- Хортой сөрөг vs Бүтээлч шүүмжлэл ялгалт}

Stage~1-ээс Сөрөг гэж ангилагдсан текстүүд Stage~2 руу дамжина. Энэ шатанд
хоёр ангилалт (binary classification) хийнэ:

\begin{itemize}
  \item \textbf{Хортой сөрөг:} Доромжлол, хувь хүнд чиглэсэн довтолгоо, сөрөг
        үгс, кибер дарамт зэрэг устгах шаардлагатай агуулга.
  \item \textbf{Бүтээлч шүүмжлэл:} Аргумент, нотолгоонд суурилсан сөрөг
        үнэлгээ, бодлогыг шүүмжлэх, асуудлыг тодорхойлох зэрэг хамгаалах
        шаардлагатай агуулга.
\end{itemize}

Stage~2-т мөн MN-BERT нарийн тохируулгатай загвар болон суурь загваруудыг (LR, NB, SVM)
хэрэгжүүлж харьцуулна. Энэ шатны тусгай бэрхшээл нь Хортой сөрөг болон Бүтээлч шүүмжлэл
хоёрын хоорондын семантик ойролцоо байдал бөгөөд хоёулаа сөрөг хандлагыг
агуулдаг боловч илэрхийлэх хэлбэр, зорилго нь ялгаатай.

Ялгалтын үндсэн хоёр шалгуур нь: (1)~\textbf{чиглэл} --- агуулга хувь хүнд
чиглэсэн эсвэл асуудалд чиглэсэн; (2)~\textbf{зорилго} --- гомдоох, дарамтлах
зорилготой эсвэл сайжруулах санал болгох зорилготой. Монгол хэлний ирони болон
давхар утгатай хэллэгт энэ хоёр шалгуур хоорондоо давхцаж болох бөгөөд энэ
нь Stage~2-ийн загварын хамгийн гол хүндрэл болдог. Энэхүү нарийн ялгааг
тодорхойлоход MN-BERT-ийн self-attention механизмын контекст мэдрэмж чухал
үүрэг гүйцэтгэнэ.

\subsection{Хоёр шатлалт загварын давхаргын архитектур}

Хоёр шатлалт загварын давхаргын бүтцийг Зураг~\ref{fig:two-stage-arch}-д
харуулав. Stage~1 болон Stage~2 тус бүр ижил MN-BERT суурь дээр суурилах
боловч тус тусдаа нарийн тохируулга хийгдсэн бөгөөд гаралтын classifier
head нь өөр өөр ангиллын тоотой.

\begin{figure}[H]
\centering
\scalebox{0.76}{%
\begin{tikzpicture}[
  layer/.style={draw, rounded corners=4pt, minimum width=7.5cm,
                minimum height=1.40cm, text width=7.0cm,
                align=center, font=\small, line width=0.7pt},
  hlayer/.style={draw, rounded corners=4pt, minimum width=3.5cm,
                 minimum height=1.40cm, text width=3.2cm,
                 align=center, font=\small, line width=0.7pt},
  arr/.style={->, >=Stealth, thick, line width=0.8pt},
]

% ── Shared MN-BERT backbone ───────────────────────────────────────────────
\node[layer, fill=gray!12, draw=gray!50] (inp) at (0, 0)
  {\textbf{Оролт токенчлол}\\[3pt]
   {\scriptsize [CLS] $t_1 \ldots t_n$ [SEP] \enspace·\enspace max\_length\,=\,128}};

\node[layer, fill=yellow!16, draw=yellow!65!black] (emb) at (0,-2.60)
  {\textbf{Embedding давхарга (768-хэмжээст)}\\[3pt]
   {\scriptsize Token Emb + Position Emb + Segment Emb}};

\node[layer, fill=blue!12, draw=blue!50] (enc) at (0,-5.20)
  {\textbf{BERT Encoder --- 12 Transformer блок}\\[3pt]
   {\scriptsize Multi-Head Self-Attention (12 толгой) + LayerNorm + Feed-Forward}};

\node[layer, fill=blue!22, draw=blue!60] (cls) at (0,-7.80)
  {\textbf{[CLS] Pooled вектор --- 768-хэмжээст}};

\node[layer, fill=orange!12, draw=orange!55] (drop) at (0,-10.40)
  {\textbf{Dropout}\enspace(p\,=\,0.1)};

% ── Stage 1 head (left) ───────────────────────────────────────────────────
\node[hlayer, fill=orange!20, draw=orange!60] (lin1) at (-2.6,-13.60)
  {\textbf{Linear}\\[3pt]{\scriptsize 768\,$\to$\,3}};

\node[hlayer, fill=green!16, draw=green!55] (out1) at (-2.6,-16.20)
  {\textbf{Softmax (Stage~1)}\\[3pt]
   {\scriptsize Эерэг · Саармаг · Сөрөг}};

% ── Stage 2 head (right) ──────────────────────────────────────────────────
\node[hlayer, fill=teal!14, draw=teal!55] (lin2) at (2.6,-13.60)
  {\textbf{Linear}\\[3pt]{\scriptsize 768\,$\to$\,2}};

\node[hlayer, fill=red!14, draw=red!50] (out2) at (2.6,-16.20)
  {\textbf{Softmax (Stage~2)}\\[3pt]
   {\scriptsize Хортой сөрөг · Бүтээлч шүүмжлэл}};

% ── Shared backbone arrows ────────────────────────────────────────────────
\draw[arr] (inp)  -- (emb);
\draw[arr] (emb)  -- (enc);
\draw[arr] (enc)  -- (cls);
\draw[arr] (cls)  -- (drop);

% Fork to two heads (0.80cm gap before branching)
\draw[arr] (drop.south) -- ++(0,-0.80) coordinate (fork)
           -- node[above, font=\scriptsize, text=green!55!black] {\textit{Stage~1}}
           (-2.6,-11.90) -- (lin1.north);
\draw[arr] (fork)
           -- node[above, font=\scriptsize, text=teal!55!black] {\textit{Stage~2}}
           (2.6,-11.90) -- (lin2.north);

\draw[arr] (lin1) -- (out1);
\draw[arr] (lin2) -- (out2);

\end{tikzpicture}%
}
\caption{Хоёр шатлалт MN-BERT загварын давхаргын архитектур. Хоёр загвар тус бүр
  ижил 12-блокт BERT encoder-ыг (768-хэмжээст hidden state) ашигладаг боловч
  тус тусдаа нарийн тохируулга (fine-tuning) хийгдсэн.
  Stage~1 нь 3-ангиллын (Linear 768$\to$3),
  Stage~2 нь 2-ангиллын (Linear 768$\to$2) ангиллын толгойтой (classifier head).}
\label{fig:two-stage-arch}
\end{figure}

%-------------------------------------------------------------------------------
\section{MN-BERT загварын нарийн тохируулга}

MN-BERT нь Монгол хэлний кирилл корпус дээр урьдчилан сургагдсан (pre-trained)
BERT загвар юм. Нарийн тохируулга процесст дараах тохиргоонуудыг ашиглана.

\subsection{Загварын гиперпараметрүүд}

\begin{table}[H]
\centering
\caption{Хоёр шатлалт (Stage~1/Stage~2) загварын гиперпараметрүүд. Эцсийн нэг
  шатлалт MN-BERT загварын тохиргоог (Focal Loss, max\_len=256 г.м.)
  Хүснэгт~\ref{tab:hyperparams-mnbert}-аас үзнэ үү.}
\label{tab:hyperparams}
\begin{tabular}{@{}ll@{}}
\toprule
\textbf{Параметр} & \textbf{Утга} \\
\midrule
Learning rate             & $2 \times 10^{-5}$ \\
Batch size                & 16 \\
Epoch тоо                 & 3--5 (early stopping) \\
Max sequence length       & 128 токен \\
Optimizer                 & AdamW \\
Weight decay              & 0.01 \\
Warmup steps              & Нийт алхмын 10\% \\
Loss function (Stage~1)   & Cross-entropy \\
Loss function (Stage~2)   & Binary cross-entropy \\
\bottomrule
\end{tabular}
\end{table}

\subsection{Нарийн тохируулга стратеги}

Нарийн тохируулга нь хоёр шатанд тус тусдаа хийгдэнэ. Stage~1-ийн загвар нь бүх
dataset (Эерэг, Саармаг, Сөрөг) дээр гурван ангиллын cross-entropy
loss-оор сургагдана. Stage~2-ийн загвар нь зөвхөн Сөрөг шошготой
дэд өгөгдлийн багц дээр binary cross-entropy loss-оор сургагдана.

Хоёр загвар хоёулаа MN-BERT-ийн ижил урьдчилсан сургалтын жинг (pre-trained
weights) эхний цэг болгон ашигладаг боловч нарийн тохируулгын дараа тус тусдаа
тохируулагдсан (specialized) загварууд үүснэ.

\textbf{Early stopping.} Validation loss 2 epoch дараалсан буурахгүй бол
сургалтыг зогсооно. Энэ нь хэт тохируулга (overfitting)-аас сэргийлэх
стандарт арга юм.

\textbf{Class-weighted loss.} Ангиллуудын тэнцвэргүй тохиолдолд loss
function-д class weight нэмж, цөөнх ангилалд (minority class) илүү өндөр
торгууль (penalty) тооцоолно. Weight нь ангилал бүрийн жишээний
урвуу давтамжаас тооцогдоно.

%-------------------------------------------------------------------------------
\section{Суурь загварууд}

MN-BERT-ийн гүйцэтгэлийг үнэлэхийн тулд гурван суурь загвартай харьцуулна.
Суурь загварууд бүгд ижил урьдчилсан боловсруулалтын дамжлагаар боловсруулагдсан
өгөгдлийн дээр сургагдана.

\begin{table}[H]
\centering
\caption{Суурь загваруудын тохиргоо}
\label{tab:baseline-config}
\begin{tabular}{@{}lll@{}}
\toprule
\textbf{Загвар} & \textbf{Шинж чанар} & \textbf{Тохиргоо} \\
\midrule
Logistic Regression & TF-IDF (unigram + bigram) & L2 регуляризаци, $C=1.0$ \\
Naive Bayes         & TF-IDF (unigram + bigram) & Multinomial NB \\
SVM                 & TF-IDF (unigram + bigram) & Linear kernel, $C=1.0$ \\
\bottomrule
\end{tabular}
\end{table}

TF-IDF (Term Frequency--Inverse Document Frequency) нь текстийн статистик
шинж чанарыг тоон вектор хэлбэрт хөрвүүлэх арга бөгөөд unigram болон bigram
хослолыг ашиглана. Эдгээр суурь загварууд нь MN-BERT-ийн нэмүү үнэ цэнийг
(value-add) тодорхойлоход туслах benchmark болно.

%-------------------------------------------------------------------------------
\section{Тэнцвэргүй өгөгдлийн шийдэл}

Бодит орчны өгөгдөлд ангиллуудын тоон тэнцвэр хангагдах нь ховор. Тухайлбал,
Хортой сөрөг агуулга нь Бүтээлч шүүмжлэл-ээс тоогоор хамаагүй олон, Саармаг нь Эерэг-ээс
давамгайлах зэрэг тэнцвэргүй байдал гарч болно. Энэ асуудлыг шийдвэрлэхэд
дараах аргуудыг хэрэгжүүлнэ:

\begin{itemize}
  \item \textbf{SMOTE (Synthetic Minority Over-sampling Technique):} Суурь
        загваруудын (LR, NB, SVM) TF-IDF шинж чанарын дээр SMOTE ашиглан
        цөөнх ангилалд хиймэл жишээ (synthetic sample) үүсгэнэ.
  \item \textbf{Focal Loss + жинлэсэн түүвэрлэлт:} Эцсийн нэг шатлалт MN-BERT
        загварт Focal Loss ($\gamma{=}2.0$)-ийг ангиллын урвуу давтамжийн жин
        $\alpha$ болон \texttt{WeightedRandomSampler}-тэй хослуулна (хоёр шатлалт
        хувилбарт жинлэсэн cross-entropy).
  \item \textbf{Stratified sampling:} Өгөгдлийн багцын хуваалт (train/val/test) дээр
        ангилал бүрийн харьцааг тэнцүү хадгална.
\end{itemize}

%-------------------------------------------------------------------------------
\section{Gradio хэрэглэгчийн интерфейс}

Системийн хэрэглэгчийн интерфейсийг Gradio framework ашиглан хэрэгжүүлнэ.
Gradio нь Python хэл дээр суурилсан, хурдан прототип бүтээхэд тохиромжтой
вэб framework юм. Интерфейс нь хоёр архитектурыг дэмждэг бөгөөд хэрэглэгч
радио товчоор аль нэгийг сонгох боломжтой. Интерфейсийн гол функцууд:

\begin{enumerate}
  \item \textbf{Текст оруулах:} Хэрэглэгч нэг текст эсвэл олон текст (batch)
        оруулж ангилуулах боломжтой.
  \item \textbf{Архитектур сонгох:} Радио товчоор \textit{Flat} (нэг шатлалт,
        4 ангилал) эсвэл \textit{Two-stage} (Stage~1 $\to$ нөхцөлт Stage~2)
        аль нэгийг сонгоно.
  \item \textbf{Ангиллын үр дүн харуулах:} Сонгосон архитектурын дагуу
        ангиллын үр дүнг итгэлцлийн хувь (confidence score)-тай хамт харуулна.
  \item \textbf{Дүрслэл (visualization):} Ангиллын тархалтын график,
        confidence distribution зэргийг дүрсэлж харуулна.
  \item \textbf{Batch шинжилгээ:} CSV файл хэлбэрээр олон текст нэг дор
        оруулж, үр дүнг татаж авах (download) боломжтой.
\end{enumerate}

\begin{figure}[H]
\centering
\scalebox{0.76}{%
\begin{tikzpicture}[
  wbox/.style={draw, rounded corners=5pt, minimum width=10.0cm,
               minimum height=1.50cm, text width=9.5cm,
               align=center, font=\small, line width=0.8pt},
  sbox/.style={draw, rounded corners=5pt, minimum width=10.0cm,
               minimum height=1.80cm, text width=9.5cm,
               align=center, font=\small, line width=0.8pt,
               fill=purple!8, draw=purple!45},
  bbox/.style={draw, rounded corners=5pt, minimum width=4.6cm,
               minimum height=1.50cm, text width=4.2cm,
               align=center, font=\small, line width=0.8pt},
  decbox/.style={diamond, draw, aspect=2.6, minimum width=3.6cm,
                 inner sep=2pt, align=center, font=\small,
                 fill=yellow!20, draw=yellow!70!black, line width=0.8pt},
  arr/.style={->, >=Stealth, thick, line width=0.9pt},
]

% ── SHARED TOP ───────────────────────────────────────────────────────────
\node[wbox, fill=blue!14, draw=blue!50] (inp) at (0, 0)
  {\textbf{Текст оруулах}\\[3pt]
   {\scriptsize Монгол кирилл · Gradio форм · Нэг эсвэл batch текст}};

\node[wbox, fill=orange!14, draw=orange!55] (pre) at (0,-3.00)
  {\textbf{Урьдчилсан боловсруулалт}\\[3pt]
   {\scriptsize Кирилл шүүлт · SentencePiece tokenization · max\_length\,=\,128}};

\node[sbox] (sel) at (0,-5.60)
  {\textbf{Архитектур сонгох (радио товч)}\\[3pt]
   {\scriptsize \textit{Flat}: нэг шатлалт, 4 ангилал
    \enspace|\enspace
    \textit{Two-stage}: Stage~1 $\to$ нөхцөлт Stage~2}};

% ── BRANCHES ─────────────────────────────────────────────────────────────
% sel.south = -5.60 - 0.90 = -6.50 ; fork 0.60 below = -7.10
% flat/s1pre north = -9.15  (center -9.90, half 0.75)
\node[bbox, fill=green!14, draw=green!55] (flat) at (-2.8,-9.90)
  {\textbf{MN-BERT (4-ангилал)}\\[3pt]
   {\scriptsize Эерэг · Саармаг · Хортой сөрөг · Бүтээлч шүүмжлэл}};

\node[bbox, fill=blue!14, draw=blue!50] (s1pre) at (2.8,-9.90)
  {\textbf{Stage~1: MN-BERT}\\[3pt]
   {\scriptsize 3-ангилал · Эерэг · Саармаг · Сөрөг}};

% Fork from selector
\coordinate (fork) at (0,-7.10);
\draw (sel.south) -- (fork);
\draw[arr] (fork)
           -- node[above, font=\scriptsize, text=green!55!black] {\textit{Flat}}
           (-2.8,-7.10) -- (flat.north);
\draw[arr] (fork)
           -- node[above, font=\scriptsize, text=blue!55!black] {\textit{Two-stage}}
           (2.8,-7.10) -- (s1pre.north);

% dec center -12.60 ; dec.north = -11.908 ; s1pre.south = -10.65 ; gap = 1.26 cm
\node[decbox] (dec) at (2.8,-12.60) {Сөрөг уу?};

% s2 center -15.30 ; s2.north = -14.55 ; dec.south = -13.292 ; gap = 1.26 cm
\node[bbox, fill=teal!12, draw=teal!50] (s2) at (2.8,-15.30)
  {\textbf{Stage~2: MN-BERT}\\[3pt]
   {\scriptsize 2-ангилал · Хортой сөрөг · Бүтээлч шүүмжлэл}};

% ── SHARED BOTTOM ────────────────────────────────────────────────────────
% out center -18.10 ; out.north = -17.35 ; s2.south = -16.05 ; gap = 1.30 cm
\node[wbox, fill=red!10, draw=red!45] (out) at (0,-18.10)
  {\textbf{Үр дүн ба итгэлцэл}\\[3pt]
   {\scriptsize Ангиллын шошго · Confidence \%}};

% viz center -20.80 ; viz.north = -20.05 ; out.south = -18.85 ; gap = 1.20 cm
\node[wbox, fill=violet!10, draw=violet!45] (viz) at (0,-20.80)
  {\textbf{Дүрслэл ба тайлан}\\[3pt]
   {\scriptsize Confidence график · Batch CSV татаж авах}};

% ── ARROWS ───────────────────────────────────────────────────────────────
\draw[arr] (inp) -- (pre);
\draw[arr] (pre) -- (sel);

% Two-stage internal
\draw[arr] (s1pre) -- (dec);
\draw[arr] (dec.south) -- node[right, font=\scriptsize, xshift=3pt] {Тийм} (s2.north);

% "Үгүй": dec.west = (1.0, -12.60) → (0, -12.60) → out.north
\draw[arr] (dec.west) -- node[above, font=\scriptsize] {Үгүй}
           (0,-12.60) -- (out.north);

% Flat → out: descend then bend right, enter box top from above at x=-1.5
\draw[arr] (flat.south) -- (-2.8,-16.50) -- (-1.5,-16.50) -- (-1.5,-17.35);

% Stage 2 → out: descend then bend left, enter box top from above at x=+1.5
\draw[arr] (s2.south) -- (2.8,-16.50) -- (1.5,-16.50) -- (1.5,-17.35);

\draw[arr] (out) -- (viz);

\end{tikzpicture}%
}
\caption{Gradio интерфейсийн ажиллагааны бүтэц: хэрэглэгч радио товчоор
  \textit{Flat} (нэг шатлалт 4-ангилал) эсвэл \textit{Two-stage}
  (Stage~1 $\to$ нөхцөлт Stage~2) архитектурыг сонгоно.
  Flat загвар нь текстийг нэг алхмаар 4 ангиллын аль нэгд хуваарилах бол,
  Two-stage загварын ``Тийм'' салаа Stage~2-оор Хортой сөрөг болон
  Бүтээлч шүүмжлэлийг нарийвчлан ялгадаг.}
\label{fig:gradio-flow}
\end{figure}

%-------------------------------------------------------------------------------
\section{Үнэлгээний арга зүй}

Загваруудын гүйцэтгэлийг дараах стандарт NLP үнэлгээний үзүүлэлтүүдээр хэмжинэ.

\subsection{Ангиллын үзүүлэлтүүд}

\begin{itemize}
  \item \textbf{Accuracy} --- нийт зөв ангилагдсан жишээний хувь.
  \item \textbf{Precision} --- тухайн ангилалд хамааруулсан жишээнүүдээс хэдэн
        хувь нь үнэхээр тэр ангилалд хамаарах.
  \item \textbf{Recall} --- тухайн ангилалд бодитоор хамаарах жишээнүүдээс хэдэн
        хувь нь зөв илрүүлэгдсэн.
  \item \textbf{$F_1$-score} --- Precision болон Recall-ийн гармоник дундаж.
        Тэнцвэргүй өгөгдлийн багцын нөхцөлд Accuracy-ээс илүү бодитой үнэлгээ өгдөг.
  \item \textbf{Confusion Matrix} --- ангилал бүрийн зөв болон буруу
        таамаглалын матриц.
  \item \textbf{Статистик ач холбогдол} --- McNemar тест болон bootstrap
        итгэлцлийн интервалаар үр дүнгийн найдвартай байдлыг нотлох.
\end{itemize}

\subsection{Харьцуулалтын арга зүй}

Загвар бүрийг ижил test set дээр үнэлнэ. Харьцуулалтыг шударгаар хийхийн тулд:
\begin{enumerate}
  \item Бүх загвар ижил train/val/test хуваалтыг ашиглана.
  \item Бүх загвар ижил урьдчилсан боловсруулалтын дамжлагыг ашиглана (MN-BERT-ийн хувьд
        tokenizer нь өөрийнх, суурь загваруудын хувьд TF-IDF).
  \item Test set нь сургалтын явцад огт ашиглагдаагүй --- зөвхөн эцсийн
        үнэлгээнд нэг удаа ашиглагдана.
  \item Macro-average $F_1$-score-ийг гол харьцуулалтын үзүүлэлт болгоно.
\end{enumerate}

%-------------------------------------------------------------------------------
\section{Технологийн стек}

Системийг хэрэгжүүлэхэд ашиглагдах гол технологиудыг хүснэгт~\ref{tab:tech-stack}-д нэгтгэв.

\begin{table}[H]
\centering
\caption{Системийн технологийн стек}
\label{tab:tech-stack}
\begin{tabular}{@{}ll@{}}
\toprule
\textbf{Зориулалт} & \textbf{Технологи} \\
\midrule
Програмчлалын хэл       & Python 3.10+ \\
Загварын framework       & PyTorch, Hugging Face Transformers \\
Урьдчилсан загвар       & MN-BERT (Hugging Face Hub) \\
Суурь загварууд          & scikit-learn \cite{pedregosa2011} \\
Тэнцвэргүй өгөгдөл      & imbalanced-learn (SMOTE) \\
Хэрэглэгчийн интерфейс  & Gradio \\
Өгөгдлийн боловсруулалт  & pandas, NumPy \\
Дүрслэл                  & matplotlib, seaborn \\
Хувилбар хяналт          & Git, GitHub \\
\bottomrule
\end{tabular}
\end{table}

%-------------------------------------------------------------------------------
\section{Бүлгийн дүгнэлт}

Энэ бүлэгт хоёр шатлалт ангиллын системийн ерөнхий архитектур, Stage~1 болон
Stage~2 загваруудын зохион байгуулалт, MN-BERT нарийн тохируулгын стратеги,
суурь загваруудын тохиргоо, тэнцвэргүй өгөгдлийн шийдлүүд, Gradio
интерфейсийн бүтэц, болон үнэлгээний арга зүйг тодорхойлсон. Дараагийн
бүлэгт эдгээр загваруудыг бодитоор сургаж, туршилтын үр дүнг танилцуулна.
